\section*{الفصل \ref{ch:g2:ice-water-steam} --- ثلج وماء وبخار}

\begin{solution}{exo:g2:ice-water-steam:1}
دلاية الثلج: صلبة. والبحر: سائلة. والماء الخفيّ فوق الفوهة:
غازية (بخار ماء). والصقيع: صلبة. والمطر: سائلة.
\end{solution}

\begin{solution}{exo:g2:ice-water-steam:2}
رجل الثلج \omterm{def:g2:ice-water-steam:meltfreeze}{ينصهر}؛ والبركة تتبخّر؛ وغبش المرآة \omterm{def:g2:ice-water-steam:evapcond}{تكاثف}؛
والغدير \omterm{def:g2:ice-water-steam:meltfreeze}{يتجمّد}.
\end{solution}

\begin{solution}{exo:g2:ice-water-steam:3}
الدفء يدفع نحو الحالة الغازية (من الصلبة إلى السائلة إلى
الغازية)؛ والبرد يردّ نحو الصلبة (من الغازية إلى السائلة إلى
الصلبة).
\end{solution}

\begin{solution}{exo:g2:ice-water-steam:4}
لا --- فالزجاج لا يرشح. بل \omterm{def:g2:ice-water-steam:evapcond}{تكاثف} بخار الماء الخفيّ الذي في هواء
الغرفة على جدار العلبة الخارجي البارد، فصار قطرات مرئية.
\end{solution}

\begin{solution}{exo:g2:ice-water-steam:5}
النَّفَس يحمل دائمًا بخار ماء لا يُرى. والهواء المتجمّد يكاثفه
في الحال قطيرات دقيقة --- أي سحابة مرئية. أما الهواء الدافئ
فيدعه متبخّرًا، فلا يبقى ما يُرى.
\end{solution}

\begin{solution}{exo:g2:ice-water-steam:6}
\omterm{def:g2:ice-water-steam:evapcond}{التبخّر}. والشمس تضيف الدفء، والدفء هو الذي يرسل الماء إلى
الهواء غازًا أسرع.
\end{solution}

\begin{solution}{exo:g2:ice-water-steam:7}
العمود حارق --- أشدّ حرارة من ماء الحمّام الساخن. والقاعدة: لا
تضع يدًا في العمود أبدًا، وارفع الأغطية بعيدًا عن وجهك، ولا
تفعل ذلك إلا وشخص كبير بجوارك.
\end{solution}

\begin{solution}{exo:g2:ice-water-steam:8}
تُدفئ الشمس البحر \omterm{def:g2:ice-water-steam:evapcond}{فيتبخّر} بعض الماء طائرًا صاعدًا من غير أن
يُرى. وعاليًا، في البرد، \omterm{def:g2:ice-water-steam:evapcond}{يتكاثف} قطيرات غيوم. ثم تثقل القطيرات
وتسقط مطرًا يجري راجعًا إلى البحر --- مستعدًّا للدورة من جديد.
\end{solution}

\begin{solution}{exo:g2:ice-water-steam:9}
في الفجوة الصافية ماء في حالته الغازية --- بخار ماء لا يُرى.
ولا \omterm{def:g2:ice-water-steam:evapcond}{يتكاثف} قطيرات بيضاء تصنع العمود إلا أعلى بقليل، حيث يبرد
الهواء. فالبخار الحقيقي لا يُرى؛ أما الأبيض الذي نسمّيه بخارًا
فهو قطرات ماء سائل دقيقة.
\end{solution}

\begin{solution}{exo:g2:ice-water-steam:10}
\omterm{def:g2:ice-water-steam:evapcond}{تبخّر} الماء في هواء الغرفة، قليلًا كل يوم. ولتبرئة القطّ: غطِّ
كأسًا مماثلة بصحن أو بغشاء بلاستيكي ودَعِ الكأسين جنبًا إلى
جنب. وبعد أسبوع تبقى المغطّاة ممتلئة بينما ينزل مستوى المكشوفة
--- فالماء يهرب إلى الهواء، ولا قطّ يشرب من وراء غطاء.
\end{solution}
